\subsection{Finite training and a changing continuation}\label{sec:r24learning}
A small gradient is useful only when its target is the future implemented by
the returned policy. We now make both the training budget and this target
condition explicit. The result concerns the trainable square-activation
critic in Section~\ref{sec:r23policy}; it neither assumes a fixed feature
regression nor asserts convergence for every neural architecture.

Fix a symmetric matrix $M$ and numbers $0<m\leq L$ with
$mI\preceq M\preceq LI$. For a possibly rectangular hidden-weight matrix
$W\in\mathbb R^{k\times d}$, $k\geq d$, put
\[
 E(W)=W'W-M,\qquad f(W)=\tfrac14\|E(W)\|_F^2,
 \qquad \nabla f(W)=WE(W).
\]
Every entry of $W$ is trainable. Consider the actual hidden-weight update
\begin{equation}
 W_{j+1}=W_j-\alpha W_jE(W_j),\qquad
 0<\alpha\leq[8(L+m)]^{-1}.
 \label{eq:r24update}
\end{equation}

\begin{proposition}[Finite-step continuation learning]\label{prop:r24learning}
If $e_0=\|E(W_0)\|_F\leq m/2$, then for every $j\geq0$,
\begin{equation}
 \|E(W_j)\|_F\leq (1-\alpha m)^{j/2}e_0,
 \qquad \sigma_{\min}(W_j)\geq\sqrt{m/2}.
 \label{eq:r24contraction}
\end{equation}
Consequently a strictly positive target $e_*$ is attained after at most
\begin{equation}
 N(e_*)=\begin{cases}
 0,&e_0\leq e_*,\\
 \left\lceil\displaystyle\frac{2\log(e_0/e_*)}
 {-\log(1-\alpha m)}\right\rceil,&e_0>e_*.
 \end{cases}
 \label{eq:r24budget}
\end{equation}
These statements do not require $M$ and $W_0'W_0$ to commute.
\end{proposition}

\begin{proof}
In the region $\|E(W)\|_F\leq m/2$, one has
$\sigma_{\min}(W)^2\geq m/2$ and $\|W\|_2^2\leq L+m/2$.
On the line segment joining $W$ and $W-\alpha WE(W)$,
$\|V\|_2\leq(3/2)\sqrt{L+m/2}$. Differentiating the gradient gives
$\|\nabla^2 f(V)\|\leq3\|V\|_2^2+L\leq8(L+m)$, where the operator
norm is induced by the Frobenius norm. The descent inequality therefore
implies $f(W-\alpha\nabla f(W))\leq f(W)-\alpha\|\nabla f(W)\|_F^2/2$.
Since $\|WE(W)\|_F^2\geq(m/2)\|E(W)\|_F^2=2mf(W)$, the loss contracts
by $1-\alpha m$. This also preserves the region, so induction proves
\eqref{eq:r24contraction}. Solving the resulting geometric inequality
proves \eqref{eq:r24budget}.
\end{proof}

The basin is a checkable sufficient condition, not an assertion about an
arbitrary random initialization. For reuse across targets, if the previous
residual is at most $e_{\rm old}$ and the target change is at most
$\Delta_M$ in Frobenius norm, then
\begin{equation}
 e_0\leq e_{\rm old}+\Delta_M.
 \label{eq:r24warm}
\end{equation}
Thus the new target's rank basin and its inner training budget can be checked
before applying \eqref{eq:r24update}. Failure of this test does not prove
that reuse fails; it means that this sufficient learning bound is unavailable.

\paragraph{Arithmetic in the training loop.}
Write an implemented step as the right side of \eqref{eq:r24update} plus
$\Delta_j$, with $\|\Delta_j\|_F\leq\omega_j$. Let
$q=\sqrt{1-\alpha m}$ and $s=\sqrt{L+m/2}$. The certified scalar recursion
\begin{equation}
 b_{j+1}=q b_j+2s\omega_j+\omega_j^2,
 \qquad b_0\geq\|E(W_0)\|_F,
 \label{eq:r24arithmetic}
\end{equation}
continues to enclose the residual as long as every $b_j\leq m/2$.
Indeed the exact step has residual at most $q b_j$ and spectral norm at
most $s$; expanding the Gram matrix after adding $\Delta_j$ gives the two
additional terms. Outward bounds are required for $m,L,b_0,q,s,\omega_j$
and the displayed arithmetic. An ideal real-arithmetic update is recovered
by setting $\omega_j=0$; zero rounding is not assumed for stored computations.

\paragraph{Training error is not policy error.}
At date $t+1$, let $S_{t+1}$ be the frozen evaluation target used for
training and let $P_{t+1}^K$ be the value matrix of the ultimately returned
policy. They need not coincide. Define
\begin{equation}
 M_{t+1}=dS_{t+1},\qquad
 \delta_{t+1}=\|S_{t+1}-P_{t+1}^K\|_F,
 \qquad a_{t+1}=b_{t+1}/d+\delta_{t+1},
 \label{eq:r24drift}
\end{equation}
where $b_{t+1}$ encloses the trained Gram residual relative to $M_{t+1}$.
At the terminal date set $a_T=0$. With $Z_t$ and $F_t$ as in
Section~\ref{sec:r23policy}, the same subtraction used in
Proposition~\ref{prop:r23factor} now gives
\begin{equation}
 \|D_t\|_F\leq \|Z_t\|_F+
 \beta\|B_t\|_2\|F_t\|_2
       \bigl(b_{t+1}/d+\delta_{t+1}\bigr).
 \label{eq:r24learningpolicy}
\end{equation}
This is a bound against the returned policy's own future. A gradient against
an earlier target cannot remove $\delta_{t+1}$.

For example, suppose a verified allowance $\bar\eta$ is assigned to the
relative policy certificate. If
\[
 z_t+c_t\delta_{t+1}<\sqrt{\bar\eta q_t r_t},\qquad
 z_t\geq\|Z_t\|_F,\quad
 c_t\geq\beta\|B_t\|_2\|F_t\|_2>0,
\]
then it suffices to train to
\begin{equation}
 b_{t+1}\leq d\left\{
 \frac{\sqrt{\bar\eta q_t r_t}-z_t}{c_t}-\delta_{t+1}\right\}.
 \label{eq:r24allocated}
\end{equation}
For $c_t=0$, only the actor residual matters. The realized actor and target
drift must still be re-evaluated before this account is used. The formula
is not a termination theorem for an arbitrary outer policy-iteration loop.
It supplies a finite inner-loop budget under explicit premises, followed
by an own-policy check. Policy-gradient convergence in structured control
provides a complementary method-level analysis
\citep{r24fazel2018}; the proposition here concerns the continuation learner,
not a transfer of that global result to a different algorithm.

The corresponding work account charges basin checks, target evaluation,
all hidden-weight updates, actor solves, failed checks, and writes. A smaller
sufficient number of updates is not by itself a total-work advantage over
a conventional representation. The matched quadratic receives the same
information and remains an essential comparator.
